% Part: first-order-logic
% Chapter: model-theory
% Section: partial-iso

\documentclass[../../../include/open-logic-section]{subfiles}

\begin{document}

\olfileid{mod}{bas}{dlo}
\section{稠密线性序}

\begin{defn}
  一个\emph{无端点稠密线性序}是指满足以下语句的结构
  $\Struct{M}$，其语言仅包含一个二元谓词符号~$<$：
  \begin{enumerate}
  \item $\lforall[x][\lnot x < x]$;
  \item $\lforall[x][\lforall[y][\lforall[z][(x < y \lif (y < z \lif x
    <z ))]]]$;
  \item $\lforall[x][\lforall[y][(x< y \lor \eq[x][y] \lor y < x)]]$;
  \item $\lforall[x][\lexists[y][x < y]]$;
  \item $\lforall[x][\lexists[y][y < x]]$;
  \item $\lforall[x][\lforall[y][(x < y \lif \lexists[z][(x < z \land
        z < y))]]]$。
 \end{enumerate}
\end{defn}

\begin{thm}\ollabel{thm:cantorQ}
  任意两个可数的无端点稠密线性序都是同构的。
\end{thm}

\begin{proof}
  设 $\Struct{M_1}$ 和 $\Struct{M_2}$ 是可数的无端点稠密线性序，
  其论域分别为 ${<_1} = \Assign{<}{M_1}$ 和 ${<_2} =
  \Assign{<}{M_2}$，
  并令 $\PIso{I}$ 为它们之间所有部分同构的集合。
  $\PIso{I}$ 非空，因为至少 $\emptyset \in \PIso{I}$。
  我们证明 $\PIso{I}$ 满足往返性质。那么 $\Struct{M_1} \iso[p] \Struct{M_2}$，
  从而根据 \olref[pis]{thm:p-isom1}，定理得证。

  为证明 $\PIso{I}$ 满足向前性质，设 $p \in \PIso{I}$ 且对 $i = 1$,
  \dots,~$n$ 有 $p(a_i) = b_i$，
  不妨设 $a_1 <_1 a_2 <_1 \cdots <_1
  a_n$。
  给定 $a \in \Domain{M_1}$，按如下方式寻找 $b \in \Domain{M_2}$：
  \begin{enumerate}
  \item 若 $a <_1 a_1$，则取 $b \in \Domain{M_2}$ 使得 $b <_2 b_1$；
  \item 若 $a_n <_1 a$，则取 $b \in \Domain{M_2}$ 使得 $b_n <_2 b$；
  \item 若对某个 $i$ 有 $a_i <_1 a <_1 a_{i+1}$，则取 $b \in \Domain{M_2}$ 使得 $b_i <_2 b <_2 b_{i+1}$。
  \end{enumerate}
  由于 $\Struct{M_2}$ 是无端点稠密线性序，总能找到满足所需性质的 $b$。
  定义 $q = p \cup \{ \langle a, b \rangle \}$，
  则 $q \in \PIso{I}$ 即为 $p$ 所需的扩张。这就证明了向前性质。
  向后性质的证明是类似的。因此 $\Struct{M_1} \iso[p]
  \Struct{M_2}$；
  由 \olref[pis]{thm:p-isom1} 可得 $\Struct{M_1} \iso
  \Struct{M_2}$。
\end{proof}

\begin{prob}
  通过验证 $\PIso{I}$ 满足向后性质，完成
  \olref[mod][bas][dlo]{thm:cantorQ} 的证明。
\end{prob}

\begin{rem}
  设 $\Struct{S}$ 是任意一个可数的无端点稠密线性序。
  那么（根据 \olref{thm:cantorQ}）$\Struct{S} \iso \Struct{Q}$，
  其中 $\Struct{Q} = (\Rat, <)$ 是以有理数集 $\Rat$ 为论域的可数稠密线性序。
  现在再次考虑 \olref[thm]{remark:R} 中的结构~$\Struct{R} = (\Real, <)$。
  我们曾看到存在一个可数结构~$\Struct{S}$ 使得 $\Struct{R} \elemequiv
  \Struct{S}$。但 $\Struct{S}$ 是一个可数的无端点稠密线性序，
  因此它同构（从而初等等价）于结构~$\Struct{Q}$。
  由初等等价的传递性，$\Struct{R} \elemequiv \Struct{Q}$。
  （我们本可以直接通过同样的往返论证建立 $\Struct{R} \iso[p] \Struct{Q}$ 来证明这一点。）
\end{rem}

\end{document}